\section{实验与结果}

He等人在ImageNet 2012分类数据集（1000类，1.2M训练图像）上进行了详尽的实验。核心对比包括：

\subsection{退化问题验证}
34层Plain网络与34层ResNet的对比表明，ResNet的Top-5错误率为12.1\%，而Plain网络仅为7.9\%——残差结构降低了4.2\%的错误率。更重要的是，Plain网络出现了退化（34层误差高于18层），而ResNet未出现退化，深度增加始终带来性能改善。

\subsection{与主流网络比较}
ResNet-152在单模型下达到4.49\%的Top-5错误率，集成6个模型后进一步降至3.03\%，刷新了当时的最佳记录。相比之下，VGG-16的Top-5错误率为6.43\%，GoogleNet为6.67\%。同时ResNet的计算复杂度（FLOPs）更低，体现了其高效性。

\subsection{CIFAR-10实验}
在CIFAR-10数据集上，ResNet-110（110层）仅约1.7M参数，达到了6.43\%的错误率，优于当时其他方法。进一步实验中，一个1202层的ResNet仍然具有竞争力（7.93\%错误率），但出现了轻微过拟合，表明深度并非无限有益，但残差结构至少使超深网络训练成为可能。

\subsection{分析实验}
论文还比较了不同快捷连接方式的影响：恒等映射（无参数）表现最佳（7.9\%误差），投影连接（$1\times1$卷积）次之（8.2\%），全零填充最差（9.2\%）。这证实了恒等映射的优越性。此外，通过分析残差块的梯度流，论文证明了残差结构可以使误差信号直接回传到浅层，有效缓解梯度消失，这是残差网络成功的关键原因之一。

整体实验表明，残差学习框架在多种数据集和网络深度下均能稳定改善训练，且不引入额外参数或显著增加计算量。